\chapter{The chygraph equation}
\label{ch:chyequation}

% Part V, added 2026-09-25 from ~/Downloads/chygraph_master_equation/draft.tex,
% Secs. 1-5. No numbers are computed here; every equation is either one the
% book already carries, re-indexed, or the general form the earlier chapters
% are special cases of.

Every chapter of Parts~II and~III runs the same recursion, and
Chapter~\ref{ch:outlook}'s Table~\ref{tab:through} lists what changes from one
to the next: what the message carries, and what is summed inside a complex.
This part writes the recursion down once, in a form general enough that every
row of that table is a substitution into it, gives its two levels the two
names they have in the literature, and then asks where else in mathematics the
same equations appear. The answer is in more places than the earlier chapters
acknowledged, and each place supplies something --- a theorem, a word, a
criterion --- that the book has been reaching for.

This chapter fixes the object. Sections~\ref{sec:ce-known}
to~\ref{sec:ce-yields} state what was known, what is added, the representation
the recursion runs on, the instance-level pair, the ensemble pair, and what
the pair yields. The five chapters after it are the connections, one per
field, each opened by a paragraph saying what the connection is and closed by
what it gives the book and what the book gives back.

\section{What was known and what is added}
\label{sec:ce-known}
\index{density evolution|textbf}
\index{belief propagation!on a chygraph|textbf}

The recursion has two levels, and the literature keeps them apart.

\emph{Chygraph belief propagation} is the instance-level pair,
Eqs.~\eqref{eq:chyup} and~\eqref{eq:chydown} below: messages on one
realisation, each updated from every incoming message but the one it answers.
Once the factor graph of Sec.~\ref{sec:ce-representation} is named, this is
belief propagation on it, and the chapter says so. What is added at this level
is not the recursion but the factor graph, which is canonical given the
chy-adjacency matrix alone, and in which every complex appears once as a
variable and once as a factor, so that the same object is iterated at every
level of nesting.

\emph{The chygraph equation} is the ensemble pair,
Eqs.~\eqref{eq:ensplus} and \eqref{eq:ensminus} below: the same recursion
averaged over a layered random chygraph, with generating functions acting
on laws and the excess bracket falling on the layer of arrival. In coding theory the ensemble version of
belief propagation is called \emph{density evolution}, after
\citet{richardson2008}, and the chygraph equation is the density evolution of chygraph belief
propagation. This is the level at which the book's thresholds, amplitudes and
free energies are computed, and the one the book can claim.

One precedent to acknowledge before claiming it. Multi-edge-type ensembles of
low-density parity-check codes \citep[Ch.~7]{richardson2008} already carry
layer-indexed degree distributions, node-perspective and edge-perspective
generating functions, and a density-evolution recursion in which the excess
falls on the edge type of arrival; their node-perspective and edge-perspective
distinction is the distinction between $\Phi$ and $\bar\Phi$, so the excess
bracket already exists there in one direction. That is the two-level chygraph ---
variables inside checks --- with a parity kernel and a log-likelihood message.
What the chygraph equation adds relative to it is the general kernel; the
nesting, with a complex's own state and its derivative at threshold
(Sec.~\ref{sec:ce-yields}); the compound complexes on which component and
cardinality differ; and the physics readout.

\section{The representation}
\label{sec:ce-representation}
\index{factor graph!of a chygraph|textbf}

Three commitments, each implicit somewhere in the book and here made at once.

\begin{model}{The factor graph of a chygraph}
\begin{rules}
\rl{1}{State} Every vertex is a variable. Atom or complex, vertex $i$
carries a state $\sigma_{i}\in\Omega_{i}$.
\rl{2}{Kernel} Every complex is also a factor. Complex $a$ carries a
non-negative function $\chi_{a}(\sigma_{a};\sigma_{\del a})$ of its own state
and of the states of its members $\del a$.
\rl{3}{Wiring} The factor graph is $\alpha+I$. Its variable nodes are all
vertices, its factor nodes are the complexes, and factor $a$ is joined to
variable $a$ and to the variables in $\del a$.
\end{rules}
\end{model}

What the rules produce is one bipartite graph, drawn in
Fig.~\ref{fig:factorgraph}, on which the recursion of the next section is
ordinary belief propagation, whatever the depth of nesting. For Part~III,
where no complex has a container, the complex variables are pendant and
the graph is the ordinary factor graph of the hypergraph model; the
construction earns its keep only when complexes contain complexes.


Rule~1 is where the models differ. For an Ising atom $\Omega_{i}=\{\pm1\}$.
For a nested clause of Sec.~\ref{sec:nested}, $\Omega_{a}$ is the clause's
truth value. For a hypergraph interaction with no state of its own,
$\Omega_{a}$ is a single point and everything below collapses to the two steps
of Chapter~\ref{ch:statmech}. For percolation, $\Omega_{a}$ records whether
$a$ is in the giant component, which is a state of every complex, atom or not.

Rule~2 is where the Hamiltonian lives, and nowhere else. An atom has
$\del i=\emptyset$ and its kernel is its own weight,
$\chi_{i}(\sigma_{i})=e^{\beta B\sigma_{i}}$ for the Ising field $B$ of
Eq.~\eqref{eq:Hising}. The distinction Chapter~\ref{ch:chygraphs} drew
between cardinality $c$ and intra-complex component $s$ lives inside
$\chi_{a}$: if the kernel factorises over the components of the interior
hypergraph, the message a member receives only ever sees that member's
component, which is why the intra-complex generating functions $G$ of
Sec.~\ref{sec:ensembles} count components and not members.

Rule~3 is the statement the book has been making in words since
Sec.~\ref{sec:treelike}. With $\alpha_{ij}=1$ when $i\in\del j$, the
bipartite graph of rule~3 has incidence matrix $\alpha+I_{C}$, the identity on
the complexes; \emph{treelike over complexes} is the statement that this
factor graph is locally a tree. Every inclusion $i\in a$ carries two messages,
both functions of $\sigma_{i}$,
\begin{align*}
  h_{i\to a}&\colon \Omega_{i}\to\mathbb{R}_{+}
  &&\text{member to container; percolation's } Q^{ai}_{+},\\
  u_{a\to i}&\colon \Omega_{i}\to\mathbb{R}_{+}
  &&\text{container to member; percolation's } Q^{ia}_{-},
\end{align*}
each defined up to a positive factor. The letters are
Chapter~\ref{ch:statmech}'s, and for an Ising variable the two functions are
the two exponentials $e^{h\sigma}$ and $e^{u\sigma}$ whose exponents that
chapter iterates.


\section{Chygraph belief propagation}
\label{sec:ce-bp}
\index{interior operator|textbf}
\index{interior sum}
\index{exterior product|textbf}

On the factor graph of rule~3 the sum-product update is one equation per
direction of an inclusion:
\begin{align}
  h_{i\to a}(\sigma_{i}) &=
    \prod_{b\ni i,\; b\ne a} u_{b\to i}(\sigma_{i})\;
    \sum_{\sigma_{\del i}} \chi_{i}(\sigma_{i};\sigma_{\del i})
    \prod_{j\in\del i} h_{j\to i}(\sigma_{j}),
  \label{eq:chyup}\\
  u_{a\to i}(\sigma_{i}) &=
    \sum_{\sigma_{a}}\ \sum_{\sigma_{\del a\setminus i}}
    \chi_{a}(\sigma_{a};\sigma_{\del a})
    \prod_{j\in\del a,\; j\ne i} h_{j\to a}(\sigma_{j})
    \prod_{b\ni a} u_{b\to a}(\sigma_{a}),
  \label{eq:chydown}
\end{align}
with the one-vertex marginal read off as
\begin{equation}
  b_{i}(\sigma_{i})\ \propto\ \prod_{b\ni i} u_{b\to i}(\sigma_{i})
  \sum_{\sigma_{\del i}}\chi_{i}(\sigma_{i};\sigma_{\del i})
  \prod_{j\in\del i} h_{j\to i}(\sigma_{j}).
  \label{eq:chyread}
\end{equation}
The pair is belief propagation on the factor graph of $\alpha+I$, and is
exact iff that factor graph is a forest, which is iff the complex family is
\emph{Berge-acyclic}: \citet{fagin1983} define Berge-acyclicity as the
acyclicity of the incidence graph, so the second ``iff'' is a definition
and the first is the tree exactness of belief propagation. It is the
forest condition of Sec.~\ref{sec:mergeerror}. That is a stronger condition than the three names of
Chapter~\ref{ch:exactness}, which are $\alpha$-acyclicity and belong to
the region-graph calculation of Chapter~\ref{ch:gbp}; two triangles
sharing an edge separate the two, and Chapter~\ref{ch:mobius} says how.

\begin{figure}[tb]
\centering
\begin{adjustbox}{max width=0.96\linewidth}
\begin{tikzpicture}[x=1.15cm,y=1.05cm]
% ---- left: the inclusion diagram (Sec. 2.7 convention) ----
\begin{scope}
  \node[atm] (a1) at (0,0) {};
  \node[atm] (a2) at (0.8,0) {};
  \node[atm] (a3) at (1.6,0) {};
  \node[atm] (a4) at (2.6,0) {};
  \node[cpx] (A)  at (0.8,1.1) {};
  \node[cpx] (Bc) at (1.7,2.2) {};
  \draw[inc] (a1)--(A) (a2)--(A) (a3)--(A);
  \draw[inc] (A)--(Bc) (a4)--(Bc);
  \node[lb,below=1pt of a1] {$1$};
  \node[lb,below=1pt of a2] {$2$};
  \node[lb,below=1pt of a3] {$3$};
  \node[lb,below=1pt of a4] {$4$};
  \node[lb,left=2pt of A] {$a$};
  \node[lb,left=2pt of Bc] {$b$};
  \node[tb,anchor=north] at (1.3,-0.55) {inclusion diagram};
\end{scope}
% ---- right: the factor graph alpha + I ----
\begin{scope}[xshift=5.4cm]
  % variables (all vertices)
  \node[atm] (v1) at (0,0) {};
  \node[atm] (v2) at (0.8,0) {};
  \node[atm] (v3) at (1.6,0) {};
  \node[atm] (v4) at (3.4,0) {};
  \node[cpx] (vA) at (1.6,1.7) {};
  \node[cpx] (vB) at (2.8,3.6) {};
  % factors (one per complex)
  \node[cpxc] (fA) at (0.8,0.9) {};
  \node[cpxc] (fB) at (2.6,2.7) {};
  \draw[inc] (v1)--(fA) (v2)--(fA) (v3)--(fA) (fA)--(vA);
  \draw[inc] (vA)--(fB) (v4)--(fB) (fB)--(vB);
  \node[lb,below=1pt of v1] {$1$};
  \node[lb,below=1pt of v2] {$2$};
  \node[lb,below=1pt of v3] {$3$};
  \node[lb,below=1pt of v4] {$4$};
  \node[lb,left=2pt of vA] {$\sigma_{a}$};
  \node[lb,right=2pt of vB] {$\sigma_{b}$};
  \node[lb,left=2pt of fA] {$\chi_{a}$};
  \node[lb,right=2pt of fB] {$\chi_{b}$};
  % the two messages on the inclusion a in b, as parallel arrows offset
  % to either side of the edge from sigma_a to chi_b
  \draw[msg] ($(vA)!0.18!(fB)+(0.16,-0.16)$) -- ($(vA)!0.82!(fB)+(0.16,-0.16)$)
     node[lb,pos=0.55,below right=-1pt] {$h_{a\to b}$};
  \draw[msg] ($(vA)!0.82!(fB)+(-0.16,0.16)$) -- ($(vA)!0.18!(fB)+(-0.16,0.16)$)
     node[lb,pos=0.45,above left=-1pt] {$u_{b\to a}$};
  \node[tb,anchor=north] at (1.7,-0.55) {factor graph of $\alpha+I$};
\end{scope}
\end{tikzpicture}
\end{adjustbox}
\caption{\textbf{The representation.} Left: a chygraph with a triangle $a$ on
atoms $1,2,3$ and a complex $b$ containing $a$ and atom $4$, drawn as an
inclusion diagram in the convention of Sec.~\ref{sec:drawing}. Right: the
factor graph of rule~3. Every vertex, atom or complex, is a variable node
(circles and open squares); every complex is in addition a factor node (filled
squares), joined to its own variable and to its members. The complex $a$
therefore appears twice, once as a variable carrying $\sigma_{a}$ and once as
the kernel $\chi_{a}$, and the inclusion $a\in b$ carries the two messages of
Sec.~\ref{sec:ce-representation}. Treelike over complexes is the statement
that this bipartite graph is locally a tree. The wiring is the chy-adjacency
matrix with the identity added on the complexes, and nothing else.}
\label{fig:factorgraph}
\end{figure}

The empty cases are the ones the book has used so far. For an atom the sum in
Eq.~\eqref{eq:chyup} is the single term $\chi_{i}(\sigma_{i})$, and
Eq.~\eqref{eq:chyup} is the up step, Eq.~\eqref{eq:up}, in exponential form:
a product of functions of one Ising variable is a sum of fields. For a complex
nothing includes, the last product in Eq.~\eqref{eq:chydown} is empty; if in
addition $\Omega_{a}$ is a point, the sum over $\sigma_{a}$ disappears and
Eq.~\eqref{eq:chydown} is the emitted field, Eq.~\eqref{eq:emitted}, whose
$\bar Z(\sigma_{0})$ is the right-hand side evaluated at the two values of
$\sigma_{i}$. Everything in Part~III is one or other of these two cases, and
the two are one equation.

\paragraph{Operator form.} The same pair written without indices separates
what is structure from what is model. Let $\otimes$ be the pointwise product
of functions on $\Omega_{i}$ --- the \emph{exterior} product, addition of
fields, a relay with no model in it --- and let $\Tint_{a}$ be the
\emph{interior} operator, the contraction of $\chi_{a}$ against the messages
on all of its legs but one. Then
\begin{equation}
  \begin{aligned}
  h_{i\to a}
    &=\Bigl(\bigotimes_{b\ni i,\,b\ne a} u_{b\to i}\Bigr)\otimes
      \Tint_{i}\bigl[h_{\cdot\to i}\bigr],\\
  u_{a\to i}
    &=\Tint^{(i)}_{a}\Bigl[h_{\cdot\to a}\,;\ \bigotimes_{b\ni a}u_{b\to a}\Bigr],
  \end{aligned}
  \label{eq:chyop}
\end{equation}
the semicolon separating the member legs from the complex's own leg, and the
superscript $(i)$ naming the leg left open. With nesting the operator has
$c-1$ member legs and one own leg, the latter present only where
$\Omega_{a}$ is not a point; Chapter~\ref{ch:operads} counts arities this
way. The Hamiltonian enters $\Tint$ and
nowhere else. Chapter~\ref{ch:localization} is the one place in the book
where $\Tint$ is not an enumeration --- there it is a Schur complement,
Eq.~\eqref{eq:schur} --- and Eq.~\eqref{eq:chyop} is why nothing else in that
chapter had to change.

\paragraph{Term by term against the percolation pair.} Percolation was solved
first because it is the case in which every one of these objects is a
number, and Table~\ref{tab:ce-termbyterm} reads the two pairs against each
other.

\begin{table}[tb]
\centering\small
\begin{adjustbox}{max width=\linewidth}
\begin{tabular}{lll}
\hline\hline
in Eqs.~\eqref{eq:chyup}--\eqref{eq:chydown}
  & in Eqs.~\eqref{eq:Qminus}--\eqref{eq:Qplus} & role\\
\hline
$\prod_{b\ni i,\,b\ne a}$ & $[\bar\Phi^{l}_{k}]_{k=m}$
  & other containers, excess on the arrival layer\\
$\sum_{\sigma_{\del i}}\chi_{i}\prod_{j\in\del i}$ & $G^{l}_{k}$
  & all members\\
$\prod_{b\ni a}$ & $\Phi^{l}_{k}$
  & all containers\\
$\sum_{\sigma_{\del a\setminus i}}\chi_{a}\prod_{j\ne i}$
  & $[\bar G^{l}_{k}]_{k=m}$
  & other members, excess on the arrival layer\\
$\sum_{\sigma_{a}}$ & implicit
  & the complex's own state\\
\hline\hline
\end{tabular}
\end{adjustbox}
\caption{\textbf{Term by term.} Each factor of the instance-level pair against
the generating function that averages it in the percolation map of
Chapter~\ref{ch:percolation}. The two excess brackets fall on the two sums
that exclude the leg being answered, one over containers and one over
members; the two ordinary functions fall on the two that do not. The last row
is the one piece with no counterpart in Chapter~\ref{ch:statmech}.
Percolation has it implicitly, because being in the giant component is a
state of every complex; nested satisfiability, Sec.~\ref{sec:nested}, needs
it explicitly, and its $\delta$ and $\gamma$ are that sum.}
\label{tab:ce-termbyterm}
\end{table}

\section{The chygraph equation}
\label{sec:ce-equation}
\index{chygraph equation|textbf}
\index{generating function!acting on a law|textbf}

Section~\ref{sec:substitution} made one substitution to pass from a
realisation to an ensemble: the argument of a generating function becomes a
measure, and the product becomes a convolution, Eq.~\eqref{eq:conv}. That
covered the up step and left the down step in words. Here is the same
substitution stated once for both.

Let a generating function act on a law $P$ over messages through a symmetric
multilinear map $T=(T_{n})_{n\ge0}$ by pushing forward product measures:
\begin{equation}
  F\act{T}P\;=\;\sum_{n}f_{n}\,(T_{n})\push\,P^{\otimes n},
  \qquad F(x)=\sum_{n}f_{n}x^{n}.
  \label{eq:act}
\end{equation}
With $T=\otimes$ this is the convolution of Eq.~\eqref{eq:conv},
$F\conv P=\sum_{n}f_{n}P^{\ast n}$, since the law of a product of independent
messages is the convolution of their laws. With $T=\Tint_{l}$ it is the down
step: draw $n$ independent fields from $P$, push them through the interior of
a layer-$l$ complex, and weight by the probability that the complex offers
$n$ legs. Both steps are therefore a generating function acting through a
multilinear map, and the asymmetry between them is only which map.

Write $P^{ml}_{+}$ for the law of the message a layer-$l$ complex sends to a
layer-$m$ container and $P^{ml}_{-}$ for the law of the message it sends to a
layer-$m$ member, $2L^{2}$ laws in all, which is Chapter~\ref{ch:percolation}'s
index set with a law in place of a number. Write $P^{l\cdot}_{+}$ for the
family over the source layer. Then, with the member count in the second line
drawn from the joint intra-complex generating function $\bar G^{l}$ over
target layers, excess on layer $m$,
\begin{align}
  P^{ml}_{+} &= \Bigl(\mathop{\conv}_{k}\,[\bar\Phi^{l}_{k}]_{k=m}\conv P^{lk}_{-}\Bigr)
               \conv \Bigl(G^{l}\act{\Tint_{l}} P^{l\cdot}_{+}\Bigr),
  \label{eq:ensplus}\\
  P^{ml}_{-} &= [\bar G^{l}]_{k=m}\act{\Tint^{(m)}_{l}}
               \Bigl(P^{l\cdot}_{+}\,;\ \mathop{\conv}_{k}\Phi^{l}_{k}\conv P^{lk}_{-}\Bigr),
  \label{eq:ensminus}
\end{align}
with the readout, no excess anywhere,
\begin{equation}
  P^{l}=\Bigl(\mathop{\conv}_{k}\Phi^{l}_{k}\conv P^{lk}_{-}\Bigr)\conv
        \Bigl(G^{l}\act{\Tint_{l}}P^{l\cdot}_{+}\Bigr).
  \label{eq:ensread}
\end{equation}
The pair is the density evolution of chygraph belief propagation over a
layered random chygraph, and it is what this part calls the \emph{chygraph
equation}. Each line is one line of Eq.~\eqref{eq:chyop} with every product
over a random number of legs replaced by the action Eq.~\eqref{eq:act}, and
the bracket $[\,\cdot\,]_{k=m}$ is Fig.~\ref{fig:themap}'s: excess on the
layer arrived through, ordinary on every other.

\paragraph{Substitutions.} Table~\ref{tab:ce-substitutions} lists what
$\otimes$ and $\Tint_{l}$ are in each of the book's calculations. In the first
row reachability factorises, so $\Tint_{l}$ is itself a product,
$(T_{n})\push\,\delta_{Q}^{\otimes n}=\delta_{Q^{n}}$, and
Eqs.~\eqref{eq:ensplus}--\eqref{eq:ensminus} reduce to compositions of
generating functions: they \emph{are} Eqs.~\eqref{eq:Qplus}--\eqref{eq:Qminus}.
That is the precise sense in which percolation is the case that compresses,
and the sense in which Chapter~\ref{ch:potts}'s $q\to1$ limit lands on a
generating function. In the third row, with one atom layer and no complex
having a container, Eq.~\eqref{eq:ensplus} is the ensemble map
Eq.~\eqref{eq:ensemblemap} and Eq.~\eqref{eq:ensminus} is the sentence
Sec.~\ref{sec:substitution} left in words: ``$Q^{lk}$ is the law of the field
a complex emits when its other members carry fields drawn from $P$''. The
fourth row is Sec.~\ref{sec:replicas-chy}: the replica saddle point iterates
a function on the $2^{n}$ replica states, and its interior sum is
Eq.~\eqref{eq:Xichy}. The fifth is Sec.~\ref{sec:onestep}, where the
substitution is made a second time and the law is a law over laws. The last
row is the one that needs the sum over $\sigma_{a}$.

\begin{table}[tb]
\centering\small
\begin{adjustbox}{max width=\linewidth}
\begin{tabular}{L{0.95in}L{1.25in}L{0.85in}L{1.55in}}
\hline\hline
model & object iterated & $\otimes$ & $\Tint_{l}$\\
\hline
percolation (Chs.~\ref{ch:percolation}--\ref{ch:epidemics})
  & a number, $P=\delta_{Q}$ & $\times$ & $\times$ over the component\\
Ising, one instance (Sec.~\ref{sec:cavity})
  & a field $h$, i.e.\ $\nu_{h}(\sigma)$ & field addition
  & the emitted field $u(h_{1},\dots,h_{c-1})$\\
replica-symmetric ensemble (Sec.~\ref{sec:substitution})
  & a law over fields & convolution & pushforward of the emitted field\\
replica saddle point (Sec.~\ref{sec:replicas-chy})
  & a function on $2^{n}$ replica states & pointwise product
  & $\Xi_{l}$, the interior sum in replica space\\
one-step RSB (Sec.~\ref{sec:onestep})
  & a law over laws, a survey & convolution of surveys
  & pushforward, reweighted by $e^{-m\Delta F}$\\
nested satisfiability (Sec.~\ref{sec:nested})
  & a warning on a clause's truth value & AND of warnings
  & one forbidden assignment on $(\sigma_{a};\text{literals})$\\
\hline\hline
\end{tabular}
\end{adjustbox}
\caption{\textbf{Substitutions.} What the exterior product and the interior
operator are in each of the book's calculations. Read downwards the second
column is Table~\ref{tab:through}'s second column again --- a number, a field,
a law, a function on replica states, a law over laws --- and the third and
fourth columns are the two things Table~\ref{tab:through} said change from
row to row, now named as the two arguments of Eq.~\eqref{eq:act}.}
\label{tab:ce-substitutions}
\end{table}

\section{What the pair yields}
\label{sec:ce-yields}
\index{threshold tensor}
\index{counting numbers}

Three things, each of which an earlier chapter obtained for one row of
Table~\ref{tab:ce-substitutions} and which the pair gives for every row at
once.

\paragraph{The threshold.} Linearising
Eqs.~\eqref{eq:ensplus}--\eqref{eq:ensminus} about the trivial fixed point,
$\otimes$ contributes a factor one per relay and $\Tint_{l}$ contributes one
derivative per leg. There are two kinds of leg, so two numbers per layer:
\begin{equation}
  u'_{l}=\frac{\partial u}{\partial h_{j}}\Big|_{\mathrm{triv}},
  \qquad
  \hat u'_{l}=\frac{\partial u}{\partial\hat h}\Big|_{\mathrm{triv}},
  \label{eq:ownleg}
\end{equation}
the first on a member leg --- it is the $u'$ of Eq.~\eqref{eq:uprime} --- and
the second on the complex's own leg, $\hat h$ being the field on
$\sigma_{a}$.
The resulting linear operator acts on the $2L^{2}$ unknowns of the
percolation tensor $A$ of Sec.~\ref{sec:threshold}, with the elementwise
reweighting of Eq.~\eqref{eq:reweight}, $\ave{s}\to\ave{s\,u'}$ and
$\ave{\sbar}\to\ave{\sbar\,u'}$, on the member channel and a new weight
$\hat u'_{l}$ on the channel that passes through a complex's own state. When
no layer has both a state and a container, $\hat u'$ never enters and the
$2L^{2}$ operator reduces to the $L\times L$ branching matrix of
Eq.~\eqref{eq:branch}, which is why Part~III could work with the smaller
index set. The three cases in the book are the three values the own-leg
derivative can take. Percolation is $\hat u'\equiv1$: reachability passes
through a complex's own state unchanged, so every entry of the tensor is a
plain moment. Part~III is $\hat u'$ absent. Nested satisfiability is
$\hat u'\in(0,1)$ --- the $\delta$ of Eq.~\eqref{eq:nestedwp} is a
warning passing through a clause's own truth value with a weight --- and it
is the reason the percolation index set, not the $L\times L$ one, is the
general form.

\paragraph{The readout.} The classical pair on a graph is one message
equation and one marginal. On a chygraph the message equation splits in two
by direction and the marginal stays one: Eq.~\eqref{eq:chyread} on an
instance, Eq.~\eqref{eq:ensread} on the ensemble, Eq.~\eqref{eq:Sl} for
percolation, where it is $S_{l}=1-P^{l}$.

\paragraph{The functional.} Equations~\eqref{eq:chyup}--\eqref{eq:chydown}
are the stationarity conditions of the Bethe free energy on the factor graph
of $\alpha+I$. Its counting numbers are $+1$ for each factor and $1-\deg(i)$
for each variable; a complex $a$ sits in its own factor and in $\kappa_{a}$
container factors, so its variable term carries $-\kappa_{a}$, an atom's
carries $1-\kappa_{i}$, and the total is one per complex, one per atom, minus
one per inclusion. These are the weights Eq.~\eqref{eq:bethe} carries and
Eq.~\eqref{eq:Fchy} derived from the replica side, and the region graph of
Chapter~\ref{ch:gbp} produces in the treelike case, now from a third route.
Chapter~\ref{ch:mobius} says what they are.

\paragraph{What is not in the pair.} Two things the book does that the pair
does not do for it. The pair is exact on a tree and silent off one; Part~IV
and Chapter~\ref{ch:exactness} are about where the factor graph of rule~3
stops being a tree, and Chapter~\ref{ch:mobius} gives that failure its
algebraic name. And the pair carries one law per directed inclusion type,
which is replica symmetry; where that is wrong the substitution has to be
made again, and Chapter~\ref{ch:rde} says what the assumption is called in
probability and what test decides it.
